\chapter{Irritability}

\section*{Definition}
An excessive or disproportionate response to stimuli characterised by anger, frustration, impatience, and a lowered threshold for annoyance, often with verbal or behavioural outbursts.

\section*{Pathophysiology}
Irritability in perimenopause is driven by oestrogen fluctuation affecting neurotransmitter systems involved in impulse control and emotional regulation. Oestrogen modulates serotonin synthesis and receptor sensitivity in the prefrontal cortex, which governs executive control over emotions. Declining oestrogen reduces serotonergic tone, decreasing the brain's ability to inhibit anger and frustration. Progesterone fluctuations affect GABAergic inhibition, further lowering the irritability threshold. Sleep disruption (common in menopause) amplifies irritability through prefrontal cortex dysfunction.

\section*{Causes}
\begin{itemize}
  \item Perimenopausal and menopausal hormonal changes
  \item Sleep deprivation or insomnia
  \item Premenstrual syndrome or PMDD
  \item Major depressive disorder or bipolar disorder
  \item Generalised anxiety disorder
  \item Chronic stress and burnout
  \item Hyperthyroidism (thyrotoxicosis)
  \item Chronic pain conditions
  \item Substance use or withdrawal (alcohol, caffeine, nicotine)
  \item Medication side effects (corticosteroids, stimulants)
\end{itemize}

\section*{When to See a Doctor}
See your doctor if:
\begin{itemize}
  \item Severe or worsening symptoms
  \item Irritability with mood swings, depression, anxiety, sleep disturbance, or relationship difficulties, especially if affecting daily function
  \item Accompanied by other concerning signs
\end{itemize}

\section*{Related Symptoms}
\begin{itemize}
  \item Mood swings
  \item Depression
  \item Anxiety
  \item Insomnia
  \item Fatigue
\end{itemize}

\textbf{Important:} If this symptom persists, your doctor will check for serious underlying causes.

\section*{Self-Care / Home Management}
\begin{itemize}
  \item Rest and avoid activities that make it worse.
  \item Maintain adequate hydration and a balanced diet.
  \item Over-the-counter remedies may provide symptomatic relief (consult a pharmacist or doctor).
  \item Monitor symptoms and seek professional advice if they persist or worsen.
  \item Avoid self-medicating with prescription medications without medical guidance.
\end{itemize}

\section*{How Your Doctor Will Evaluate You}
\begin{itemize}
  \item A thorough history and physical examination are the first steps in evaluation.
  \item Laboratory tests (blood work, urinalysis) may be ordered based on clinical suspicion.
  \item Imaging studies (X-ray, ultrasound, CT, MRI) may be indicated when structural pathology is suspected.
  \item Specialist referral is arranged when the diagnosis remains uncertain or requires specific expertise.
\end{itemize}

\section*{Treatment Options}
\begin{itemize}
  \item Treatment targets the underlying cause identified through diagnostic evaluation.
  \item Supportive care and symptom management are provided while awaiting definitive therapy.
  \item Medications, lifestyle modifications, and physical therapies may be prescribed as appropriate.
  \item Follow-up ensures treatment effectiveness and allows timely adjustments to the care plan.
\end{itemize}

\clearpage